\documentclass[11pt,a4paper]{article}
\usepackage[T1]{fontenc}
\usepackage[utf8]{inputenc}
\usepackage[british]{babel}
\usepackage{amsmath,amsthm}
\usepackage{newtxtext,newtxmath}
\usepackage[margin=27mm]{geometry}
\usepackage{microtype}
\usepackage{booktabs,tabularx}
\usepackage{enumitem}
\usepackage[round,authoryear]{natbib}
\usepackage{xcolor}
\definecolor{linkblue}{HTML}{174A72}
\usepackage[colorlinks=true,allcolors=linkblue,pdftitle={Exercise-policy loss under computational and financial constraints: Mathematical framework}]{hyperref}
\usepackage{cleveref}
\setlength{\parindent}{1.2em}
\setlength{\parskip}{0.2em}
\setlength{\emergencystretch}{2em}
\linespread{1.06}
\newtheorem{proposition}{Proposition}
\newtheorem{assumption}{Assumption}
\newcommand{\E}{\mathbb E}
\newcommand{\Q}{\mathbb Q}
\newcommand{\F}{\mathcal F}
\newcommand{\ind}{\mathbf 1}
\newcommand{\loss}{\ell}
\newcommand{\positive}[1]{\left(#1\right)^+}
\title{Exercise-policy loss under computational\newline and financial constraints\\[0.45em]\large Mathematical framework}
\author{}
\date{26 September 2026}
\begin{document}
\maketitle
\begin{abstract}
An approximate exercise policy can be difficult to assess because its decisions are suboptimal, because the upper bound used to assess them is loose, or because the financial comparison assigns different exercise opportunities. This paper develops a framework that separates these objects in a finite-calendar stopping problem with deterministic funding and discrete cash events. A martingale decomposition distinguishes expected policy loss from dual slack. Analytically centred piecewise-linear values provide a practical upper construction. Nested calendars reduce a prescribed financial and expiry sensitivity family to its maximal endpoint for unchanged exercise instructions. A conditional terminal adjustment supplies an expected policy gain that can be isolated with a common upper payoff. Independent evaluation, simultaneous scenario limits and explicit timing accounting distinguish training error, sampling uncertainty, numerical approximation and financial specification. This methods-only version presents the model, derivations and experimental design. It reports no market observations, fitted empirical quantities or numerical research findings.
\end{abstract}
\smallskip
\noindent\textbf{Keywords:} optimal stopping; American options; martingale duality; Monte Carlo; discrete dividends; computational allocation.

\section{Research question and attribution}

The economic value of an exercise algorithm is the expected reward obtained by its stopping rule. Regression fit, numerical resolution and the number of simulated paths are means of constructing or evaluating that rule; none is itself a measure of value. A useful comparison must identify the financial problem, specify a feasible policy and state what its benchmark establishes. A finely resolved price can be a valuable numerical reference without being a known optimum. A valid upper bound can remain too loose to establish that policy loss is economically small.

The framework asks three related questions. How much expected value does a fixed exercise instruction forgo within a specified model? How much computation is needed to measure that loss with useful precision? How does the answer change when the financial inputs or opportunity set change? These questions require different comparisons. Changing the martingale used for evaluation changes bound tightness. Changing a policy changes its expected reward. Changing the exercise calendar can change the optimum even when the original policy is left untouched.

The methodological foundations are established. \citet{longstaff2001} construct exercise policies through backward continuation regression. \citet{haugh2004} develop martingale upper bounds and the use of approximate value processes. The lognormal hinge expectation used below follows the analytical pricing foundation of \citet{black1973}. The derivations here apply these foundations to an explicitly funded cash-event model and organise the resulting policy comparisons. The dual inequality, conditional Jensen argument and finite-family inference are not presented as new general theorems.

This document is the mathematical component of a wider research project. Its scope is deliberately independent of any particular observed contract, fitted parameter, sample size or empirical result. It defines the evidence that a numerical and financial investigation would need, rather than asserting that the evidence has been obtained.

\section{The funded finite-calendar stopping problem}

\subsection{State, payments and information}

Work on a filtered probability space $(\Omega,\F,(\F_t)_{t\geq0},\Q)$ carrying a Brownian motion. Let $D_0(u)>0$ denote a deterministic discount factor for a payment at date $u$. If a quoted share price $S_t$ settles at $u(t)$, define
\begin{equation}
Y_t=f_tS_t,\qquad f_t=D_0(u(t)).
\label{eq:ownership}
\end{equation}
The variable $Y$ represents discounted economic ownership under the maintained settlement convention. This transformation requires the share and exercise transfer to represent matching economic entitlements. It does not establish a financing strategy or identify the correct funding curve from market prices alone.

Between prescribed cash events, assume
\begin{equation}
 dY_t=\sigma Y_t\,dW_t^{\Q},\qquad \sigma\geq0.
\label{eq:diffusion}
\end{equation}
At an event with proposed discounted cash $c\geq0$, the detached entitlement and remaining state are
\begin{equation}
 C=\min(Y_{t-},c),\qquad Y_t=\positive{Y_{t-}-c},
 \qquad Y_t+C=Y_{t-}.
\label{eq:cash}
\end{equation}
The cap preserves non-negativity and cash accounting. This bounded cash specification is related to the stock-dependent dividend framework discussed by \citet{vellekoop2006}; the settlement and ownership equations here specify the particular maintained model. The ex-entitlement state is not a martingale across the cash detachment; between events it is. Conditional transition calculations must respect this distinction. Proposed cash, volatility and funding are model inputs, not identified expectations merely because they permit a price fit.

For a put with strike $K>0$, let $\mathcal E=\{t_1,\ldots,t_J\}$, with $0<t_1<\cdots<t_J$, be a finite set of admissible exercise rights. Valuation is at $t_0=0$. Define
\begin{equation}
 H_j=\positive{Kf_{t_j}-Y_{t_j}}.
\label{eq:reward}
\end{equation}
Write $Y_j=Y_{t_j}$ and $M_j=M_{t_j}$ when processes are indexed by the decision grid; $H_{t_j}=H_j$ identifies the reward at a calendar time. Thus $H_\tau$ denotes the reward at the selected exercise right. The exercise-settlement factor can differ from the premium-payment factor $b>0$. Values below are in discounted currency unless divided by $b$. The definition of $H_j$ presumes that the specified exercise settlement transfers the intended ownership state and entitlements. A change in deliverable or exercise funding can require a different reward.

\begin{assumption}[Admissibility]
The decision set is finite and fixed before evaluation. Rewards are adapted and integrable. Policies are stopping times taking values in $\mathcal E$. Evaluation martingales start at zero and are integrable at all rights. Whenever sampling variances are reported, the corresponding random variables have finite second moments. Learned coefficients and approximate value functions are fixed before independent evaluation paths are drawn.
\label{ass:admissible}
\end{assumption}

With deterministic initial state and model inputs, the optimal and policy values are
\begin{equation}
 V^*=\sup_{\tau\in\mathcal T(\mathcal E)}\E[H_\tau],
 \qquad V_\pi=\E[H_{\tau_\pi}],
 \qquad \loss_\pi=\frac{V^*-V_\pi}{b}\geq0.
\label{eq:loss}
\end{equation}
Cash-event observations needed for propagation do not automatically grant exercise rights. The order of an exercise decision, detachment and settlement is part of the problem definition. A finite exercise calendar represents a Bermudan problem; its numerical refinement alone does not verify the contractual exercise opportunities of an American option.

\subsection{Learned exercise instructions}

Backward regression estimates continuation values on a training sample and converts them into stopping predicates \citep{longstaff2001}. All coefficients, basis choices, positivity safeguards and tie rules form part of the frozen policy. Independent forward evaluation estimates its value without treating its training fit as realised out-of-sample accuracy. Increasing basis dimension need not improve a finite-sample rule.

When a financial scenario $\theta$ changes the ownership factor, preserve the instruction in quoted-price coordinates. If $a_j^\theta$ is the scenario share-settlement factor, the quoted state is $Q_j=Y_j^\theta/a_j^\theta$. A source predicate expressed in its original ownership coordinate receives
\begin{equation}
 T_j^\theta(Y_j^\theta)
 =Y_j^\theta\frac{a_j^0}{a_j^\theta}.
\label{eq:transfer}
\end{equation}
Its original coefficients and decision conventions are retained, but the resulting stopping time is valued with scenario rewards. This preserves the instruction for an identical history of quoted prices. It does not preserve the distribution of stopping dates when the scenario dynamics change.

\section{Loss measurement and dual slack}

\subsection{A pathwise gap}

For an admissible martingale $M$ with $M_0=0$, set
\begin{equation}
 U(M)=\max_{1\leq j\leq J}\{H_j-M_j\},\qquad
 L_\pi(M)=H_{\tau_\pi}-M_{\tau_\pi},\qquad
 G_\pi(M)=U(M)-L_\pi(M).
\label{eq:gap}
\end{equation}
Here the maximum ranges over the indexed exercise rights. The martingale dual approach is due to the published framework of \citet{haugh2004}.

\begin{proposition}[What a gap measures]
Under Assumption~\ref{ass:admissible}, $G_\pi(M)\geq0$ pathwise and
\begin{equation}
 \E[G_\pi(M)]
 =\underbrace{V^*-V_\pi}_{\text{policy loss}}
 +\underbrace{\E[U(M)]-V^*}_{\text{dual slack}}.
\label{eq:decomposition}
\end{equation}
Both terms on the right are non-negative.
\end{proposition}
\begin{proof}
The maximum defining $U$ includes the policy's stopping right, proving the pathwise inequality. Finite-horizon optional sampling gives $\E[M_\tau]=0$ for each admissible stopping time. Hence $\E[L_\pi]=V_\pi$. Also $U\geq H_\tau-M_\tau$ for every such $\tau$, so $\E[U]\geq\sup_\tau\E[H_\tau]=V^*$. Subtracting $V_\pi$ and adding and subtracting $V^*$ gives \eqref{eq:decomposition}.
\end{proof}

For a fixed policy and financial problem, a change in expected gap caused only by replacing $M$ is a change in dual slack. It is not a change in policy value. A numerical PDE value supplies an additional comparison but is not inserted as an exact $V^*$ without a justified error enclosure. Mesh agreement is evidence about numerical behaviour, not a proof of the financial model.

If two policies use the same upper payoff and martingale, then
\begin{equation}
 G_\pi(M)-G_\rho(M)=L_\rho(M)-L_\pi(M),\qquad
 \E[G_\pi-G_\rho]=V_\rho-V_\pi.
\label{eq:common-upper}
\end{equation}
The upper payoff cancels pathwise. Paired gap differences can therefore identify expected policy gains without confounding them with changes in the bound.

\subsection{Analytically centred approximate values}

Let $v_{j+1}$ be a deterministic approximation to a subsequent value function, or a function fixed by an independent training procedure. Define
\begin{equation}
 \Delta M_{j+1}=v_{j+1}(Y_{j+1})
 -\E[v_{j+1}(Y_{j+1})\mid\F_{t_j}].
\label{eq:centred}
\end{equation}
When the conditional expectation is evaluated for the actual retained function and transition, these increments have conditional mean zero in exact arithmetic. Approximation quality affects tightness rather than supplying the reason for centring. The construction is the martingale component of an approximate value process described by \citet{haugh2004}.

I represent approximate values on the non-negative state domain as
\begin{equation}
 v(y)=\alpha+\beta y+\sum_{r=1}^{R}\gamma_r\positive{y-k_r},
 \qquad k_r\geq0.
\label{eq:hinges}
\end{equation}
Over a no-cash interval of length $h$, $X=y\exp(-\sigma^2h/2+\sigma\sqrt h\,Z)$, with $Z$ standard normal. For $y,k>0$ and $\sigma^2h>0$,
\begin{align}
 \E[\positive{X-k}]&=y\Phi(d_1)-k\Phi(d_2),\\
 d_1&=\frac{\log(y/k)+\sigma^2h/2}{\sigma\sqrt h},\qquad
 d_2=d_1-\sigma\sqrt h.
\label{eq:hinge-integral}
\end{align}
This is the lognormal option integral \citep{black1973}. Zero state, strike, variance or time uses its continuous limiting form. Each term in \eqref{eq:hinges} can consequently be integrated over the full support, including the retained tail rule.

For a capped cash jump with size $c$, composition is explicit:
\begin{equation}
 v(\positive{x-c})
 =\alpha+\beta\positive{x-c}
 +\sum_{r=1}^{R}\gamma_r\positive{x-(c+k_r)}.
\label{eq:composition}
\end{equation}
The identity uses $k_r\geq0$ and includes the constant segment for $x\leq c$. Cash-phase ordering must be preserved when multiple propagation intervals are joined. An approximate quadrature expectation or truncated integration range cannot simply be substituted into \eqref{eq:centred} while retaining an exact martingale claim.

Floating-point evaluation, function compression and approximation of the dynamic programme are distinct numerical errors. Validation should check analytical limiting cases, cash accounting, tail treatment and independent pricing references. A compression tolerance for one function is not an accumulated error bound for the stopping problem.

\section{Changed opportunities and changed policies}

\subsection{Nested calendars with unchanged instructions}

Let $\theta\in\Theta$ index a prescribed finite financial family and let $\mathcal E(d)$ be nested exercise calendars. Write $V^*_\theta(d)$ for the optimum and suppose that the original policy always stops on the smallest calendar. Its reward definition and pre-terminal dynamics remain unchanged when $d$ varies within a financial scenario.

\begin{proposition}[Maximal-calendar reduction]
If $\mathcal E(d_1)\subseteq\mathcal E(d_2)$ whenever $d_1\leq d_2$, and $b_\theta>0$ is constant across these calendars, then
\begin{equation}
 \max_{\theta\in\Theta,\,d\in D}
 \frac{V^*_\theta(d)-V_{\pi,\theta}}{b_\theta}
 =\max_{\theta\in\Theta}
 \frac{V^*_\theta(d_{\max})-V_{\pi,\theta}}{b_\theta}.
\label{eq:maximal}
\end{equation}
\end{proposition}
\begin{proof}
Every stopping rule on a smaller calendar is feasible on a larger one, so the optimum is nondecreasing. The original policy value is constant within the scenario because its decisions and rewards remain unchanged. Subtracting this constant and dividing by a positive fixed premium factor preserves monotonicity. Taking the maxima gives \eqref{eq:maximal}.
\end{proof}

The statement concerns true expected losses. Sample upper limits need not increase with the calendar because sampling noise and dual slack can change. Evaluating the prespecified maximal calendar supplies a bound for smaller nested calendars; it does not claim to have priced or simulated every intermediate case. Recalibration across endpoints, a changed cash transfer or a policy allowed to use new observations would require a different comparison.

\subsection{A feasible terminal adjustment}

Suppose the original endpoint is $T$, the extended endpoint is $T+h$, and between them $Y$ is an integrable martingale with no cash event and a constant exercise threshold $k$. Define
\begin{equation}
 g(y)=\positive{k-y},\qquad
 g_h(y)=\E[g(Y_{T+h})\mid Y_T=y],\qquad
 q_h(y)=g_h(y)-g(y).
\label{eq:wait}
\end{equation}
Conditional Jensen gives $q_h(y)\geq0$. Under the lognormal transition, $g_h$ follows the corresponding analytical put integral. The endpoint condition is stronger than calendar nesting: it provides a tractable terminal replacement and dominance of earlier exercise within the added interval.

Let $A=\{\tau_\pi=T\}\in\F_T$. Change only the final instruction: if $A$ occurs, wait until $T+h$; otherwise retain the original stopping time. Denote this adjusted policy by $\widetilde\pi$.

\begin{proposition}[Expected gain from the final instruction]
The adjusted rule is feasible on the expanded calendar and
\begin{equation}
 V_{\widetilde\pi}-V_\pi=\E[\ind_Aq_h(Y_T)]\geq0.
\label{eq:terminal-gain}
\end{equation}
The conditionally integrated gain has variance no greater than the direct realised payoff difference, when the latter has finite variance.
\end{proposition}
\begin{proof}
The decision to wait is made using the $\F_T$-measurable event $A$, so the rule is a stopping time. Outside $A$, payoffs agree. On $A$, the tower property gives the expected difference as $q_h(Y_T)$. Conditional Jensen proves its sign. The integrated gain is the conditional expectation of the realised difference given $\F_T$; the law of total variance proves the variance comparison.
\end{proof}

The realised later payoff can still be smaller. Equation~\eqref{eq:terminal-gain} establishes an expected improvement under the maintained model, not a pathwise profit or executable market opportunity.

For the expanded problem, a common upper payoff isolates this gain as in \eqref{eq:common-upper}. The residual discounted loss satisfies
\begin{equation}
 V_h^*-V_{\widetilde\pi}
 =(V_0^*-V_\pi)+(V_h^*-V_0^*)-\E[\ind_Aq_h(Y_T)].
\label{eq:remaining}
\end{equation}
The policy gain need not be bounded above by the increase in the optimum: a suboptimal original rule may survive where an optimal original rule would already stop. Residual loss for a changing policy need not be monotone in the endpoint.

\subsection{Retraining as a separate comparison}

Replacing the terminal training response by $g_h$ can also change earlier continuation regressions. To isolate that effect, use the same training paths, basis, regression safeguards and financial inputs for original and revised fits, and give both policies the same final waiting instruction. Freeze the fits before fresh paired evaluation.

Equation~\eqref{eq:terminal-gain} does not imply that finite-sample retraining improves value. Changed stopping decisions alter the responses entering earlier regressions, and estimated continuation can move in either direction. A common upper payoff identifies the incremental expected value through \eqref{eq:common-upper}. Unfavourable and inconclusive paired outcomes must remain in the comparison.

\section{Conditional inference and scenario families}

\subsection{What the probability statement conditions on}

With the financial model, fitted policy and approximate values fixed, simulate independent outer paths. Let $X_i=G_\pi^{(i)}(M)/b$, with sample mean $\overline X$ and sample standard deviation $s_X$. A commonly used approximate upper limit is
\begin{equation}
 B_\pi=\overline X+t_{1-\alpha,n-1}\frac{s_X}{\sqrt n}.
\label{eq:pointwise}
\end{equation}
Since $\loss_\pi\leq\E[X_i]$, a valid upper limit for the expected gap is also an upper limit for policy loss. Student calibration is generally approximate for non-normal gaps; it is not a finite-sample guarantee. Finite variance, independent evaluation and sufficiently informative Monte Carlo samples remain necessary. More outer paths improve precision about the expected gap but do not remove its dual slack.

A conditional interval does not integrate over the uncertainty of policy training, financial calibration, contract conventions or sample selection. Training repetitions must be retained separately if their variation is investigated. Choosing a policy on the same evaluation outcomes used to advertise its bound can invalidate the selected-policy interpretation. Such a selection requires a separate evaluation design or an appropriate simultaneous guarantee.

\subsection{Coverage for a prescribed financial family}

For $m=|\Theta|$, compute a per-scenario limit at the maximal calendar using error allocation $\alpha/m$:
\begin{equation}
 B_\theta=\overline X_\theta+
 t_{1-\alpha/m,n_\theta-1}\frac{s_\theta}{\sqrt{n_\theta}}.
\label{eq:family}
\end{equation}
If each marginal limit has its stated coverage, the union bound and \eqref{eq:maximal} imply an event of probability at least $1-\alpha$ on which every prescribed scenario and smaller calendar is covered. With approximate marginal calibration, the joint statement inherits that approximation. Independence across scenarios is unnecessary for the union bound.

There is no extra multiplicity factor for calendars already dominated by the maximal endpoint. There is also no automatic extension to all policies, a continuous funding region or scenarios introduced after inspection. A missing required scenario prevents a complete-family conclusion; the family size must not shrink opportunistically. Financial quote compatibility is a different condition and must be reported separately from mathematical completeness of the conditional family.

\subsection{Pairing, independent runs and gap interactions}

Within a common-path experiment, I estimate policy gains or scenario effects from pathwise differences. For $D_i=X_i-Y_i$,
\begin{equation}
 \widehat{\operatorname{Var}}(\overline D)=\frac{s_D^2}{n}
 =\frac{s_X^2+s_Y^2-2s_{XY}}{n}.
\label{eq:paired}
\end{equation}
Discarding covariance would answer a different uncertainty question. Conversely, unrelated evaluation streams are not paired merely because policies or model labels coincide.

For independent mean estimates $R$ and $L$, possibly with different path counts,
\begin{equation}
 \widehat v=\frac{s_R^2}{n_R}+\frac{s_L^2}{n_L},\qquad
 \nu=\frac{\widehat v^2}
 {(s_R^2/n_R)^2/(n_R-1)+(s_L^2/n_L)^2/(n_L-1)}.
\label{eq:welch}
\end{equation}
The usual Welch interval uses $t_{1-\alpha/2,\nu}\sqrt{\widehat v}$. Degenerate zero-variance cases require explicit treatment. A difference between reported confidence endpoints is a descriptive comparison, not itself a confidence interval for a policy-value difference.

A financial/calendar interaction can be written in terms of expected quoted gaps as
\begin{equation}
 I_\theta=\{g_\theta(d_{\max})-g_0(d_{\max})\}
          -\{g_\theta(0)-g_0(0)\}.
\label{eq:interaction}
\end{equation}
When each financial effect is paired internally but the two calendar experiments are independent, add the variances of their paired mean effects. By \eqref{eq:decomposition}, $I_\theta$ combines interactions in true loss and dual slack. It does not identify either component alone and has no causal interpretation. Repeated dependent pointwise comparisons do not constitute a simultaneous test; intervals containing zero do not establish equivalence or exact additivity.

\section{Computational expenditure as part of the design}

\subsection{Specify the workload before comparing methods}

A time comparison is meaningful only after defining the computational unit. For a collection of frozen policies under one financial model, approximate-value construction and the common martingale can be shared. Charging this construction once per policy overstates its cost; omitting it understates the cost. A single-policy workload can have a different ranking from a repeated multi-policy workload.

The timed unit should state whether it includes simulation, deterministic approximation, conditional integration, martingale evaluation, policy realisation, estimator construction and numerical checks. Previous training, imports, data loading and file output should be reported separately if excluded. Warm-up and allocation pilots are additional expenses, not free inputs to an end-to-end valuation. Record elapsed and CPU time and the relevant hardware and thread configuration.

\subsection{Allocate paths before inspecting evaluation outcomes}

One possible allocation design uses median pilot times $t_1,t_2$ at two path counts $n_1<n_2$:
\begin{equation}
 b_T=\frac{t_2-t_1}{n_2-n_1},\qquad
 a_T=\max(0,t_1-b_Tn_1),\qquad
 \widehat T(n)=a_T+b_Tn.
\label{eq:timing}
\end{equation}
For target $B$ and path increment $q$, choose
\begin{equation}
 n(B)=q\left\lfloor\frac{B-a_T}{b_Tq}\right\rfloor,
\label{eq:allocation}
\end{equation}
subject to prescribed limits. A non-positive fitted slope or infeasible minimum is a design failure. The unclipped intercept, timing dispersion and any cap must be retained. Fix all allocations before drawing the independent evaluation streams, and do not stop a run because a confidence bound becomes favourable.

The affine model predicts time; it does not enforce it. Extrapolation beyond the pilot range, memory effects and host variability can create substantial discrepancies. Compare actual time with target time and retain overruns. A strict equal-cost ranking requires adequate timing calibration or a separately justified comparison on an observed cost frontier. Nominally equal targets alone do not establish it.

\subsection{Approximation and sampling serve different purposes}

A more detailed approximate value can tighten the dual expectation while making each outer path more expensive. Under a time constraint this can reduce the available sample size and increase the sampling radius in \eqref{eq:pointwise}. Report the mean gap, its standard error, the upper limit and actual runtime separately. A finer deterministic grid is not automatically the best use of a limited evaluation budget.

For independent-method evaluations, use \eqref{eq:welch}; for common-path evaluations, retain \eqref{eq:paired}. Report all retained policies and repetitions. Descriptive averages or medians across related cases are not independent observations of a population effect. Neither a tighter dual construction nor a larger outer sample changes a frozen policy's expected reward. Comparing these expenditures with extra policy training is a distinct resource-allocation problem.

\section{Financial validation and interpretation}

A numerical stopping problem can be well defined while its market interpretation is unresolved. Empirical application requires dated contract terms, a justified ownership state, cash information available at the valuation cutoff, and funding appropriate to the payment obligations. Event and source timestamps may describe different objects. Simultaneous records can require an explicit ordering rule; absent or stale inputs must not silently become valid observations.

Joint call and put compatibility under a shared state, cash specification, funding convention and volatility is an informative diagnostic. Separate unconstrained calibrations can remove that restriction. Compatibility with bid--ask intervals is still not identification of the correct model or an executable arbitrage. A model that fails a renewed quote diagnostic after a sensitivity change can remain a conditional mathematical case, but must not be reclassified as an accepted empirical specification because its policy loss is small.

A quoted spread offers a monetary scale. It is not automatically the cost of exercise, hedging or implementation. An upper loss limit below the spread establishes only that conditional scale comparison. A limit above it leaves the inequality unestablished; it does not prove that true loss exceeds the spread.

A prospective evaluation should freeze its observations, admissibility rule, calibration and inference before outcomes can affect method choice. If no observation passes the prescribed screen, a performance estimate is unavailable. Wider post-sample thresholds can inform a future design but cannot silently replace the primary sample. Numerical replication, admissible measurement and independent economic validation are different requirements.

\section{Conclusions and data availability}

I distinguish policy quality from the tightness of its evaluation bound and from the financial opportunities used as the benchmark. Analytically centred values can improve the measurement of loss without changing the policy. A larger calendar changes the comparison problem, while a feasible terminal adjustment changes the rule itself. Common-upper comparisons with valid martingales identify expected policy gains; finite-family bounds and explicit timing accounting delimit what the resulting evidence can support.

I use this separation to distinguish two allocation questions. Additional computation spent on the dual construction can improve knowledge of a fixed policy's loss; additional computation spent on training can change the policy's reward. Neither expenditure can be ranked solely by the resulting dual gap because that gap combines loss and slack. A comparison of training and evaluation budgets therefore requires a common financial problem, a declared computational workload and independent evaluation of the resulting policies. The framework specifies that comparison but does not identify an optimal allocation without evidence on both margins.

Predictable compensation identifies a separate source of error if the conditional mean-zero requirement is relaxed. A zero apparent gap or a common upper payoff cannot certify policy quality in its presence. I therefore retain analytical centring as a validity condition, separate from the approximation choices used to tighten the bound.

These distinctions provide a basis for a conditional computational investigation. They do not establish which specification performs best in observed markets or how the results generalise to additional state variables and uncertain funding or cash. Such conclusions require their own numerical and empirical evidence.

In this methods-only publication, I set out the model and apply the cited methods. The finite-tree example is stipulated for analytical verification. I include no raw or cleaned subscription observations, fitted market inputs, empirical summaries, research-result tables or data-derived figures. I retain the complete numerical paper separately. The LaTeX sources and build instructions for this version are available in the accompanying repository.


\bibliographystyle{plainnat}
\bibliography{references}
\end{document}
